\documentclass[10pt,a4paper]{amsart}
\usepackage[margin=19mm]{geometry}
\usepackage{amsmath,amssymb,mathtools}
\usepackage[scr=rsfs,cal=euler]{mathalfa}
\usepackage{xfrac}
\usepackage[unicode,hidelinks,bookmarks=false]{hyperref}
\newcommand{\ie}{i.e.\ }
\newcommand{\cf}{cf.\ }
\newcommand{\resp}{respectively\ }
\newcommand{\Subsection}{\subsection}
\providecommand{\nobreakdash}{\nobreak\hbox{-}}
\setlength{\emergencystretch}{2em}
\multlinegap0pt
\usepackage{bm}
\usepackage{bbm}
\usepackage{dsfont}
\usepackage{xspace}
\usepackage{cancel}
\usepackage{mathtools}
\usepackage{amsthm}
\makeatletter
\@ifundefined{theorem}{\newtheorem{theorem}{Theorem}}{}
\makeatother
\newcommand {\PPP}{\mathcal{P}}
\newcommand {\SSS}{\mathcal{S}^1}
\newcommand {\nneg}{\bm{\neg}}
\newcommand{\eqdef}{\stackrel{\mathrm{def}}{=}}
\title[Isoperimetric problem for quasi-Finsler metrics and shapes o]{Isoperimetric problem for quasi-Finsler metrics and shapes of least aerodynamic resistance}
\author{Lev Lokutsievskiy, Alexander Plakhov}
\date{Source: arXiv:2608.17783v2 (CC BY 4.0)}
\begin{document}
\maketitle

\noindent\emph{Source and licence.} Isoperimetric problem for quasi-Finsler metrics and shapes of least aerodynamic resistance. Version arXiv:2608.17783v2 (CC BY 4.0), distributed under the Creative Commons Attribution 4.0 International licence, as recorded in the arXiv licence metadata for this exact version and shown on the abstract page https://arxiv.org/abs/2608.17783v2. Full rights notice: licenses/arxiv-2608.17783-CC-BY-4.0.txt.\par\smallskip

\noindent\textit{Editorial scope.} Counted unit: the Theorem whose source label is \texttt{t\_cond} in the article (source0.tex lines 809--817, source label \texttt{t\_cond}), delivered together with the article's own complete original proof of it (source0.tex lines 820--910, one \texttt{proof} environment of 91 lines). No mathematics is written, repaired or re-derived: statement and proof are the article's own text line for line. Required context retained uncounted: S5P (lines 766--768); P+P'' (lines 770--772); PP'' (lines 774--776). Every definition, display and label of those lines is retained unchanged. Omitted from this package and not redistributed: 1--765, 769--769, 773--773, 777--808, 818--819, 911--1333. Nothing inside a retained excerpt is omitted or altered: every retained body is delivered exactly as the article's lines are written, so no omission receipt of any kind accompanies this package. Citations: the retained excerpts carry no citation command at all, so no bibliography entry of the article is transcribed and no reference is redistributed. Reference adaptation: none was needed and none was made; every reference inside the delivered unit resolves inside it, so no unresolved reference or citation is produced. Printed slips kept literal: no printed word, symbol, hypothesis or number of the counted statement or of its proof was altered. Layout and class: the article's own class is \texttt{article} with packages including pstricks, graphicx, psfrag, amsmath, amssymb, amsthm, color, bm, xcolor, hyperref, geometry; the wrapper is the lane's pinned \texttt{amsart} editorial wrapper at 10pt on A4, which restates the theorem-like environments the retained text needs and adds the article's own macro definitions that the retained text needs (\texttt{\textbackslash PPP}, \texttt{\textbackslash SSS}, \texttt{\textbackslash eqdef}, \texttt{\textbackslash nneg}), with extra wrapper packages (bm, bbm, dsfont, xspace, cancel, mathtools). The delivered numbering is the unit's own. Assets: no retained span contains a figure environment, a graphics-inclusion command, a PDF-page inclusion, a table or a TikZ picture; no image file is copied into this package, the package declares no asset dependency and no image file is redistributed with it. Licence: version arXiv:2608.17783v2 of this article is distributed under the Creative Commons Attribution 4.0 International licence, as recorded in the arXiv licence metadata for this exact version and shown on the abstract page https://arxiv.org/abs/2608.17783v2. A full rights notice is delivered at licenses/arxiv-2608.17783-CC-BY-4.0.txt. Characters: every retained excerpt is pure ASCII, so the delivered engine, XeLaTeX with its default Latin Modern text font, reports no missing character.
\begin{equation}\label{S5P}
\PPP(\theta)\, = \, \dfrac{1}{2}\int_{-\pi/2}^{\pi/2} \cos^3 \varphi\, f(\theta-\varphi)\, d\varphi,
\end{equation}

\begin{equation}\label{P+P''}
\PPP(\theta) + \PPP''(\theta)\, =\, - \int_{-\pi/2}^{\pi/2} \cos 3\varphi\, f(\theta-\varphi)\, d\varphi
\end{equation}

\begin{equation}\label{PP''}
\big(\PPP(\theta) + \PPP''(\theta) \big)'\, =\, 3\int_{-\pi/2}^{\pi/2} \sin 3\varphi\, f(\theta-\varphi)\, d\varphi.
\end{equation}

\begin{theorem}\label{t_cond}
(i)\ If $K \le 0$ and $k \ge \pi/2$ (i.e., $\nneg$C1 and $\nneg$C2), then $\PPP>0$ on $\SSS \setminus \{\pi\}$ and $\PPP+\PPP''>0$ on $\hat\SSS$.

(ii)\ If $K > 0$ and $k \ge \pi/2$ (i.e., C1 and $\nneg$C2), then $\PPP>0$ on $\SSS \setminus \{\pi\}$, and $\PPP+\PPP'' \le 0$ in $[-\theta_*,\, \theta_*]$ and $\PPP+\PPP''> 0$ on $\hat\SSS \setminus [-\theta_*,\, \theta_*]$, for some $0 < \theta_* \le \pi/6$.

(iii)\ If $K < 0$ and $k < \pi/2$ (i.e., $\nneg$C1 and C2), then $\PPP>0$ on $(-\pi/2-k,\, \pi/2+k)$ and $\PPP+\PPP'' >0$ on $\hat\SSS \cap (-\pi/2-k,\, \pi/2+k)$, and $\PPP=0$ on $\SSS \setminus (-\pi/2-k,\, \pi/2+k)$.

(iv)\ If $K \ge 0$ and $k < \pi/2$ (i.e., C1 and C2), then $\PPP>0$ on $(-\pi/2-k,\, \pi/2+k)$,\, $\PPP=0$ on $\SSS \setminus (-\pi/2-k,\, \pi/2+k)$,\, $\PPP+\PPP''\le 0$ on $[-\theta_*,\, \theta_*]$ (for some $0 < \theta_* \le \pi/6$), and $\PPP+\PPP''> 0$ on $\hat\SSS \cap (-\pi/2-k,\, \pi/2+k) \setminus [-\theta_*,\, \theta_*]$.
\end{theorem}

\begin{proof}
Note that the function $\PPP+\PPP''$ is continuous and even. The theorem follows from the following statements.
\vspace{2mm}

(a) If $k < \pi/2$ then $\PPP>0$ on the interval $(-\pi/2-k,\, \pi/2+k)$ and $\PPP=0$ outside this interval. If $k \ge \pi/2$ then $\PPP>0$ everywhere, except possibly at $\theta=\pi$ (when $k=\pi/2$).
\vspace{2mm}

(b) $\PPP + \PPP'' > 0$ on the interval $(\pi/6,\, \min\{\pi/2 + k,\, \pi \}) \setminus \{ k-\pi/6 \}$.
 %$(\pi/6,\, \pi/2 + k) \setminus \{ k-\pi/6 \}$, if $k < \pi/2$, and on $(\pi/6,\, \pi) \setminus \{ k-\pi/6 \}$, otherwise.
\vspace{2mm}

(c) The function $\dfrac{1}{\cos 3\theta}\,(\PPP(\theta) + P''(\theta))$ is non-decreasing on $[0,\, \pi/6)$. Additionally, if $k \ge \pi/2$, it is strictly increasing, and if $k < \pi/2$, it is constant on a segment $[0,\, \theta_0]$ with $\theta_0>0$ sufficiently small.
\vspace{2mm}

Let us prove these statements.\vspace{2mm}

(a) Using that $f$ is supported on $[-k,\, k]$, one sees that the integrand in \eqref{S5P} is positive when $\varphi$ is in the interval $(\theta-k,\, \theta+k) \cap (-\pi/2,\, \pi/2)$ and zero when $\varphi$ lies outside the closure of this interval. For $k < \pi/2$, this interval is nonempty, and equivalently $\PPP>0$, iff $\varphi \in (-\pi/2-k,\, \pi/2+k)$. Similarly, for $k=\pi/2$,\, $\PPP>0$ on $\SSS \setminus \{ \pi \}$, and for $k> \pi/2$,\, $\PPP>0$ everywhere.
\vspace{2mm}

(b) Take $\theta \in (\pi/6,\, \min\{\pi/2 + k,\, \pi \}) \setminus \{ k-\pi/6 \}$. Using that $\cos 3\theta$ is negative on $(-\pi/2,\, -\pi/6)$ and making the change of variable $\alpha = 3\varphi - \pi/2$, we get
\begin{equation}\label{eqn1}
\PPP(\theta) + \PPP''(\theta)\ =\
-\int_{-\pi/2}^{-\pi/6} \cos 3\varphi\, f(\theta-\varphi)\, d\varphi
-\int_{-\pi/6}^{\pi/2} \cos 3\varphi\, f(\theta-\varphi)\, d\varphi
 \end{equation}
 \begin{equation}\label{eqn2}
 \ge\
-\int_{-\pi/6}^{\pi/2} \cos 3\varphi\, f(\theta-\varphi)\, d\varphi\, =\,
\dfrac{1}{3}\int_{-\pi}^{\pi} \sin\alpha\ f\Big(\theta - \frac{\pi}{6} - \dfrac{\alpha}{3} \Big)\, d\alpha
\end{equation}
 \begin{equation}\label{eqn3}
\, =\,
\dfrac{2}{3}\int_{0}^{\pi} \sin\alpha
\Big[ f\Big(\theta - \frac{\pi}{6} - \dfrac{\alpha}{3} \Big) - f\Big(\theta - \frac{\pi}{6} + \dfrac{\alpha}{3} \Big)
\Big] d\alpha.
\end{equation}
Using that $0 \le \Big|\theta - \dfrac{\pi}{6} - \dfrac{\alpha}{3}\Big| \le \pi - \Big|\dfrac{7\pi}{6} - \theta - \dfrac{\alpha}{3}\Big| \le \pi$, we have
$$
f\Big(\theta - \dfrac{\pi}{6} - \dfrac{\alpha}{3} \Big)\ =\ f\Big(\Big|\theta - \dfrac{\pi}{6} - \dfrac{\alpha}{3} \Big|\Big)\ \ge\
f\Big(\pi - \Big|\dfrac{7\pi}{6} - \theta - \dfrac{\alpha}{3}\Big|\Big) \ =\ f\Big(\theta - \dfrac{\pi}{6} + \dfrac{\alpha}{3} \Big),$$
since $f(\beta)=f(-\beta)$ and $f(\pi-\beta)=f(\pi+\beta)$. Therefore, $\PPP(\theta) + \PPP''(\theta) \ge 0$ for all $\theta$.

Now we examine carefully the cases when $\PPP(\theta) + \PPP''(\theta) > 0$. Note that $f(\theta-\varphi)>0$ if $\varphi\in(\theta-k,\theta+k)$. If either $\theta - k<-\pi/6$ or $\theta + k > 3\pi/2$, then the first integral in the right hand site of \eqref{eqn1} is negative, and therefore, the inequality in \eqref{eqn2} is strict and hence $\PPP(\theta)+\PPP''(\theta)>0$.

The opposite case is $\theta -k \ge - \pi/6$ and $\theta + k \le 3\pi/2$. Note that $\theta-k\ne-\pi/6$ by assumptions in (b). Hence $\theta -k > - \pi/6$ and $\theta + k \le 3\pi/2$. We claim that the interval
\begin{equation}\label{intervals}
3\Big( \Big|k-\theta+\dfrac{\pi}{6}\Big|, \ \min\big\{ k+\theta-\dfrac{\pi}{6},\, \frac{13\pi}{6}-\theta-k \big\}\Big) \cap [0,\, \pi]
\end{equation}
is nonempty, and the first term in the square brackets in \eqref{eqn3} is positive for the values $\alpha$ in the interval~\eqref{intervals}, while the second term is zero, and so, the integral in \eqref{eqn3} is positive.

Indeed, one easily checks that the former interval in \eqref{intervals} is empty only for $k=\pi$, which contradicts the considered case. So the former interval in \eqref{intervals} is nonempty. From the inequalities $\theta < \pi/2+k$ (condition (b)) and $k< \theta+\pi/6$ one derives that the left endpoint of the former interval lies in $[0,\, \pi)$, and therefore, the set in \eqref{intervals} is nonempty. Each $\alpha$ in the interval \eqref{intervals} satisfies $-k+\theta-\pi/6<\alpha/3<k+\theta-\pi/6\ \implies\ \theta-\pi/6-\alpha/3 \in (-k,\, k)$, hence the first term in the square brackets in \eqref{eqn3} is positive. Finally, each $\alpha$ in the interval \eqref{intervals} satisfies $k-\theta+\pi/6 <\alpha/3< 13\pi/6-\theta-k\ \implies\ \theta-\pi/6+\alpha/3 \in (k,\, 2\pi-k)$, hence the second term in the square brackets in \eqref{eqn3} is zero.

It follows that $\PPP(\theta) + \PPP''(\theta) > 0$, except possibly for a single point $\theta = k - \pi/6$.
	\vspace{2mm}

(c) By \eqref{P+P''} and \eqref{PP''} we have
$$
I(\theta) \eqdef -\int_{-\pi/2}^{\pi/2} \sin(3\theta-3\varphi)\, f(\theta-\varphi)\, d\varphi
$$ $$
=\ -\sin 3\theta \int_{-\pi/2}^{\pi/2} \cos 3\varphi\, f(\theta-\varphi)\, d\varphi\
+\ \cos 3\theta \int_{-\pi/2}^{\pi/2} \sin 3\varphi\, f(\theta-\varphi)\, d\varphi
$$ $$ =\ \sin 3\theta\, \big( \PPP(\theta) + \PPP''(\theta) \big)\
+\ \dfrac{1}{3} \cos 3\theta\, \big( \PPP(\theta) + \PPP''(\theta) \big)'.
$$
On the other hand, making the change of variable $\psi = \theta-\varphi$ and using that $f$ is even and nonnegative, for $0 < \theta < \pi/6$ we obtain (note that function $f(\psi)\sin3\psi$ is odd)
$$
I(\theta)\ =\ -\int_{-\pi/2+\theta}^{\pi/2+\theta} \sin 3\psi\, f(\psi)\, d\psi $$ $$
=	 \int_{-\pi/2}^{-\pi/2+\theta} \sin 3\psi\, f(\psi)\, d\psi
\ -\ \int_{\pi/2}^{\pi/2+\theta} \sin 3\psi\, f(\psi)\, d\psi\ \ge\ 0.
$$
If $k \ge \pi/2$, the last inequality is strict, and if $k < \pi/2$, it becomes equality for $0\le \theta \le \pi/2-k$.

So, we have
$$
\dfrac{d}{d\theta}\, \dfrac{\PPP(\theta) + \PPP''(\theta)}{\cos 3\theta}\ =\ 3\ \dfrac{\sin 3\theta\, \big( \PPP(\theta) + \PPP''(\theta) \big)\
+\ \dfrac{1}{3} \cos 3\theta\, \big( \PPP(\theta) + \PPP''(\theta) \big)'}{\cos^2 3\theta}\ $$ $$ =\ \dfrac{3I(\theta)}{\cos^2 3\theta}\ \ge \ 0.
$$
If $k \ge \pi/2$, this inequality is strict, and if $k < \pi/2$, it becomes equality for $\theta\ge 0$ sufficiently small.
\vspace{2mm}

Let us now prove the theorem.
\vspace{2mm}

(i) We have $\PPP(0)+\PPP''(0) =-K \ge 0$. Since $k \ge \pi/2$, by (c) the function $\dfrac{1}{\cos 3\theta}(\PPP(\theta)+\PPP''(\theta))$ is strictly increasing on $(0,\,\pi/6)$. According to (b), $\PPP+\PPP'' >0$ on $(\pi/6,\,\pi) \setminus \{ k-\pi/6 \}$. Hence $\PPP+\PPP'' >0$ everywhere, except possibly at $\theta = 0,\, \pm\pi/6,\, \pi,\, \pm(k-\pi/6)$. Here and in item (ii), by (a), $\PPP>0$ everywhere, except possibly at $\theta=\pi$ (when $k=\pi/2$).

(ii) We have $\PPP(0)+\PPP''(0) =-K < 0$ and, by (c), $\dfrac{1}{\cos 3\theta}(\PPP(\theta)+\PPP''(\theta))$ is strictly increasing on $(0,\,\pi/6)$. According to (b), $\PPP+\PPP'' >0$ on $(\pi/6,\,\pi) \setminus \{ k-\pi/6 \}$. It follows that $\PPP+\PPP'' \le 0$ on a certain interval $[0,\, \theta_*]$,\, $0<\theta_*\le\pi/6$, and $\PPP+\PPP'' > 0$ on $(\theta_*,\, \pi) \setminus \{ \pi/6,\, k-\pi/6 \}$, and the same inequalities are true on the symmetric intervals $[-\theta_*,\, 0]$ and $(-\pi,\, -\theta_*) \setminus \{ -\pi/6,\, -k+\pi/6 \}$. Thus, $\PPP+\PPP'' \le 0$ on $[-\theta_*,\, \theta_*]$ and $\PPP+\PPP'' >0$ on $\hat\SSS \setminus [-\theta_*,\, \theta_*].$

(iii) We have $\PPP(0)+\PPP''(0) =-K > 0$. Since $k < \pi/2$, by (c), $\dfrac{1}{\cos 3\theta}(\PPP(\theta)+\PPP''(\theta))$ is non-decreasing on $(0,\,\pi/6)$. According to (b), $\PPP+\PPP'' >0$ on $(\pi/6,\,\pi/2+k)\setminus \{ k-\pi/6 \}$.  It follows that $\PPP+\PPP'' >0$ on $(-\pi/2-k,\, \pi/2+k) \setminus \{ \pm\pi/6,\, \pm(k-\pi/6) \}$. Here and in item (iv), by (a), $\PPP>0$ on $(-\pi/2-k,\, \pi/2+k)$ and $\PPP=0$ on $\SSS \setminus (-\pi/2-k,\, \pi/2+k)$.

(iv) We have $\PPP(0) + \PPP''(0) = -K\le 0$. Since $k < \pi/2$, by (c), $\PPP(\theta) + \PPP''(\theta) = -K\cos 3\theta \le 0$ on a certain segment $[0,\,\theta_0]$, $0<\theta_0<\pi/6$. By (c) and (b), $\PPP+\PPP'' \le 0$ on $[0,\, \theta_*]$, where $\theta_0 \le \theta_* \le \pi/6$,\, $\PPP+\PPP'' >0$ on $(\theta_*,\, \pi/2+k)\setminus \{ \pi/6,\, k-\pi/6 \}$ (and the same inequalities hold true on the symmetric segments). Thus, the last item of Theorem \ref{t_cond} is proved.
\end{proof}

\end{document}
